\documentclass[12pt]{article}
\usepackage{amsmath, amssymb, amsthm}
\usepackage{geometry}
\geometry{margin=1in}
\usepackage{fancyhdr}
% \usepackage[breaklinks=True]{hyperref}
\usepackage[bookmarks]{hyperref}
\usepackage{tikz}
\usepackage{hyperxmp}
\usepackage{enumitem}
\usetikzlibrary{calc}

\newcommand{\PP}{\mathbb{P}} % polynomials
\newcommand{\CC}{\mathbb{C}} % complex numbers
\newcommand{\RR}{\mathbb{R}} % real numbers
\newcommand{\QQ}{\mathbb{Q}} % rational numbers
\newcommand{\NN}{\mathbb{N}} % nonnegative integers
\newcommand{\calF}{\mathcal{F}}
\newcommand{\Z}[1]{\mathbb{Z}/#1\mathbb{Z}} % integers modulo k
                                            % (syntax: "\Z{k}")
\newcommand{\ZZ}{\mathbb{Z}} % integers
\newcommand{\id}{\operatorname{id}} % identity map
\newcommand{\lcm}{\operatorname{lcm}}
% Lowest common multiple. For historical reasons, LaTeX has a \gcd
% command built in, but not an \lcm command. The preceding line
% rectifies that.
\newcommand{\diag}{\operatorname{diag}} % diagonal matrix
\newcommand{\ord}{\operatorname{ord}} % order of a group element
\newcommand{\rank}{\operatorname{rank}} % rank of a matrix
\newcommand{\Hom}{\operatorname{Hom}} % group of homomorphisms
\newcommand{\End}{\operatorname{End}} % ring of endomorphisms
\newcommand{\Iso}{\operatorname{Iso}} % set of isomorphisms
\newcommand{\GL}{\operatorname{GL}} % general linear group
\newcommand{\SL}{\operatorname{SL}} % special linear group
\newcommand{\Nil}{\operatorname{Nil}} % nilradical
\newcommand{\Jac}{\operatorname{Jac}} % Jacobson radical
\newcommand{\Tor}{\operatorname{Tor}} % torsion submodule or Tor functor
\newcommand{\Ext}{\operatorname{Ext}} % Ext functor
\newcommand{\Ker}{\operatorname{Ker}} % kernel
\newcommand{\Coker}{\operatorname{Coker}} % cokernel
\newcommand{\symd}{\mathbin{\bigtriangleup}} % symmetric difference of sets
\newcommand{\powset}[2][]{\ifthenelse{\equal{#2}{}}{\mathcal{P}\left(#1\right)}{\mathcal{P}_{#1}\left(#2\right)}}
% $\powset[k]{S}$ stands for the set of all $k$-element subsets of
% $S$. The argument $k$ is optional, and if not provided, the result
% is the whole powerset of $S$.
\newcommand{\set}[1]{\left\{ #1 \right\}}
% $\set{...}$ compiles to {...} (set-brackets).
\newcommand{\seq}[1]{\left\{ #1 \right\}}
% $\seq{...}$ compiles to {...} (seq-brackets).
\newcommand{\abs}[1]{\left| #1 \right|}
% $\abs{...}$ compiles to |...| (absolute value, or size of a set).
\newcommand{\tup}[1]{\left( #1 \right)}
% $\tup{...}$ compiles to (...) (parentheses, or tuple-brackets).
\newcommand{\ive}[1]{\left[ #1 \right]}
% $\ive{...}$ compiles to [...] (Iverson bracket, aka truth value).
\newcommand{\floor}[1]{\left\lfloor #1 \right\rfloor}
% $\floor{...}$ compiles to |_..._| (floor function).
\newcommand{\underbrack}[2]{\underbrace{#1}_{\substack{#2}}}
% $\underbrack{...1}{...2}$ yields
% $\underbrace{...1}_{\substack{...2}}$. This is useful for doing
% local rewriting transformations on mathematical expressions with
% justifications. For example, try this out:
% $ \underbrack{(a+b)^2}{= a^2 + 2ab + b^2 \\ \text{(by the binomial formula)}} $
\newcommand{\horrule}[1]{\rule{\linewidth}{#1}} % Create horizontal rule command with 1 argument of height
\newcommand{\nnn}{\nonumber\\} % Don't number this line in an "align" environment, and move on to the next line.
\newcommand{\explain}[1]{\tup{\hspace{-0.5pc}\text{\begin{tabular}{c}#1\end{tabular}\hspace{-0.5pc}}}}
% $\explain{...}$ interprets the "..." as text (you can
% nest some math inside by putting it in $...$'s as you
% would in any other text block) and puts it in parentheses.

% Arrows for normal use:
\newcommand{\mono}{\hookrightarrow} % Monomorphism/injection.
\newcommand{\epi}{\twoheadrightarrow} % Epimorphism/surjection.
\newcommand{\iso}{\overset{\cong}{\to}} % Epimorphism/surjection.

% Arrows for xymatrix:
\newcommand{\arinj}{\ar@{_{(}->}} % Monomorphism.
\newcommand{\arinjrev}{\ar@{^{(}->}} % Monomorphism.
\newcommand{\arsurj}{\ar@{->>}} % Epimorphism.
\newcommand{\arelem}{\ar@{|->}} % Arrow to be used in describing a map on elements.
\newcommand{\arback}{\ar@{<-}} % Backward arrow.

%----------------------------------------------------------------------------------------
%	MAKING SUMMATION SIGNS ALWAYS PUT THEIR BOUNDS ABOVE AND BELOW
%	THE SIGN
%----------------------------------------------------------------------------------------
% The following are hacks to ensure that sums (such as
% $\sum_{k=1}^n k$) always put their bounds (i.e., the $k=1$ and the
% $n$) underneath and above the sign, as opposed to on its right.
% Same for products (\prod), set unions (\bigcup) and set
% intersections (\bigcap). Remove the 8 lines below if you do not want
% this behavior.
\let\sumnonlimits\sum
\let\prodnonlimits\prod
\let\cupnonlimits\bigcup
\let\capnonlimits\bigcap
\let\maxnonlimits\max
\let\minnonlimits\min
\let\oldlimits\lim
\renewcommand{\sum}{\sumnonlimits\limits}
\renewcommand{\prod}{\prodnonlimits\limits}
\renewcommand{\bigcup}{\cupnonlimits\limits}
\renewcommand{\bigcap}{\capnonlimits\limits}
\renewcommand{\max}{\maxnonlimits\limits}
\renewcommand{\min}{\minnonlimits\limits}
\renewcommand{\lim}{\oldlimits\limits}

%----------------------------------------------------------------------------------------
%	ENVIRONMENTS
%----------------------------------------------------------------------------------------
% The incantations below define how theorem environments
% (\begin{theorem} ... \end{theorem}) and their likes will look like.
\newtheoremstyle{plainsl}% <name>
  {8pt plus 2pt minus 4pt}% <Space above>
  {8pt plus 2pt minus 4pt}% <Space below>
  {\slshape}% <Body font>
  {0pt}% <Indent amount>
  {\bfseries}% <Theorem head font>
  {.}% <Punctuation after theorem head>
  {5pt plus 1pt minus 1pt}% <Space after theorem headi>
  {}% <Theorem head spec (can be left empty, meaning `normal')>

% Environments which make the text inside them slanted:
\theoremstyle{plainsl}
  \newtheorem{theorem}{Theorem}[section]
  \newtheorem{proposition}[theorem]{Proposition}
  \newtheorem{lemma}[theorem]{Lemma}
  \newtheorem{corollary}[theorem]{Corollary}
  \newtheorem{conjecture}[theorem]{Conjecture}
% Environments that don't:
\theoremstyle{definition}
  \newtheorem{definition}[theorem]{Definition}
  \newtheorem{example}[theorem]{Example}
  \newtheorem{exercise}[theorem]{Exercise}
  \newtheorem{examples}[theorem]{Examples}
  \newtheorem{algorithm}[theorem]{Algorithm}
  \newtheorem{question}[theorem]{Question}
 \theoremstyle{remark}
  \newtheorem{remark}[theorem]{Remark}
\newenvironment{statement}{\begin{quote}}{\end{quote}}
\newenvironment{fineprint}{\begin{small}}{\end{small}}

\newcommand{\sol}{\subsection*{Solution}}
\newcommand{\prob}[1]{\section*{Problem #1}}
\newcommand{\letterbullet}[1]{\item[\textbf{(#1)}]}
\newcommand{\compsign}{(*)}
\newcommand{\diff}[2]{\frac{d#1}{d#2}}
\newcommand{\ydiff}[1]{\diff{y}{#1}}
\newcommand{\pdiff}[2]{\frac{\partial #1}{\partial #2}}
\newcommand{\eval}[1]{\left. #1 \right\rvert}
\newcommand{\st}{\;:\;}
\newcommand{\grad}{\nabla}
\newcommand{\inv}[1]{#1^{-1}}


\newcommand{\eps}{\varepsilon}

\newcommand{\rref}[1]{\xrightarrow{#1}}
\newcommand{\interchange}{\leftrightarrow}

\newcommand{\inner}[2]{\langle #1, #2 \rangle}

\newcommand{\norm}[1]{\left\| #1 \right\|}

\newcommand{\zero}[1]{\mathbf{0}_{#1}}

\newcommand{\Real}[1]{\operatorname{Re}(#1)}

% the for all symbol pisses me off without a space after it. this fixes it
\let\oldforall\forall
\renewcommand{\forall}{\oldforall\;}

% the exists symbol pisses me off without a space after it. this fixes it
\let\oldexists\exists
\renewcommand{\exists}{\oldexists\;}

\newcommand{\recall}{\textbf{Recall:} }

\DeclareMathOperator{\proj}{proj}

\pagestyle{fancy}
\fancyhf{}
\rhead{Partial Differential Equations -- Homework 2}
\lhead{Marc DeCarlo}
\cfoot{\thepage}

\title{Partial Differential Equations -- Homework 2}
\author{Marc DeCarlo}
\date{\today}

\begin{document}

\maketitle

\prob{1}
Solve the heat flow problem with Dirichlet boundary conditions
\[
\begin{cases}
u_t = u_{xx}, & 0 < x < 2,\ t > 0, \\
u(0,t) = 0,\quad u(2,t) = 2, & t > 0, \\
u(x,0) = x + \sin(\pi x), & 0 \leq x \leq 2.
\end{cases}
\]
Identify the steady-state and transient parts of the solution.

\sol
This problem is nonhomogeneous, so the solution will contain a steady-state part and a transient part. The solution looks like the following:
\[
u(x,t) = s(x) + v(x,t),
\]
where $s(x)$ is the steady-state solution and $v(x,t)$ is the transient part of the solution. The steady-state solution is the solution to the following boundary value problem because the transient part will vanish as $t \to \infty$:
\[
\begin{cases}
s_{xx} = 0, & 0 < x < 2, \\
s(0) = 0,\quad s(2) = 2.
\end{cases}
\]
\textbf{Steady-State Solution}
\[
\begin{cases}
s_{xx} = 0, & 0 < x < 2, \\
s(0) = 0,\quad s(2) = 2.
\end{cases}
\]
To solve this, we integrate the equation twice:
\[
s_{xx} = 0 \implies s_x = C_1 \implies s = C_1 x + C_2
\]
Applying the boundary conditions, we get
\[
s(0) = C_2 = 0 \implies C_2 = 0
\]
\[
s(2) = 2C_1 = 2 \implies C_1 = 1
\]
Thus, the steady-state solution is
\[
s(x) = x
\]
\textbf{Transient Solution}
Since the solution to the entire problem is $u(x,t) = s(x) + v(x,t)$, we can substitute this into the original PDE to find the transient part of the solution:
\[
u(x,t) = s(x) + v(x,t) 
\]
We need to find $v(x,t)$ so:
\[
v(x,t) = u(x,t) - s(x) = u(x,t) - x
\]
This makes the new PDE for $v$:
\[
\begin{cases}
v_t = v_{xx}, & 0 < x < 2,\ t > 0, \\
v(0,t) = 0,\quad v(2,t) = 0, & t > 0, \\
v(x,0) = u(x,0) - s(x) = \sin(\pi x), & 0 \leq x \leq 2.
\end{cases}
\]
This is a homogeneous heat equation with homogeneous Dirichlet boundary conditions. We can solve this using separation of variables or by inspection. Let $v(x,t) = X(x)T(t)$, then we have
\[
\frac{T'}{T} = \frac{X''}{X} = -\lambda^2
\]
The $X$ part of the solution is
\[
X'' + \lambda^2 X = 0
\implies X(x) = A \sin(\lambda x) + B \cos(\lambda x)
\]
Applying the boundary conditions, we get
\[
X(0) = A \sin(0) + B \cos(0) = B = 0 \implies B = 0
\]
\[
X(2) = A \sin(2\lambda) = 0 \implies \lambda = \frac{n\pi}{2}, \ n = 1,2,3,\dots
\]
Thus, the transient solution is
\[
v(x,t) = \sum_{n=1}^{\infty} A_n \sin\left(\frac{n\pi x}{2}\right) e^{-\left(\frac{n\pi}{2}\right)^2 t}
\]
To find the coefficients $A_n$, we use the initial condition for $v$:
\[
v(x,0) = \sin(\pi x) = \sum_{n=1}^{\infty} A_n \sin\left(\frac{n\pi x}{2}\right)
\]
The only term in the series that will be nonzero is when $n=2$ because $\sin(\pi x) = \sin\left(\frac{2\pi x}{2}\right)$. Thus, we have
\[
A_2 = 1
\]
and all other $A_n$ are zero. Thus, the transient solution is
\[
v(x,t) = \sin(\pi x) e^{-\pi^2 t}
\]
\textbf{Final Solution}
Thus, the final solution to the problem is
\[
\boxed{
u(x,t) = s(x) + v(x,t) = x + \sin(\pi x) e^{-\pi^2 t}
}
\]

\prob{2}
Solve the heat flow problem with Neumann boundary conditions
\[
\begin{cases}
u_t = u_{xx}, & 0 < x < 1,\ t > 0, \\
u_x(0,t) = u_x(1,t) = 1, & t > 0, \\
u(x,0) = x + \cos^2(\pi x), & 0 \leq x \leq 1.
\end{cases}
\]
Identify the steady-state and transient parts of the solution.

\sol
This problem is nonhomogeneous, so the solution will contain a steady-state part and a transient part. The solution looks like the following:
\[
u(x,t) = s(x) + v(x,t),
\]
where $s(x)$ is the steady-state solution and $v(x,t)$ is the transient part of the solution. The steady-state solution is the solution to the following boundary value problem because the transient part will vanish as $t \to \infty$:
\[
\begin{cases}
s_{xx} = 0, & 0 < x < 1, \\
s_x(0) = 1,\quad s_x(1) = 1.
\end{cases}
\]
\textbf{Steady-State Solution}
\[
\begin{cases}
s_{xx} = 0, & 0 < x < 1, \\
s_x(0) = 1,\quad s_x(1) = 1.
\end{cases}
\]
To solve this, we integrate the equation twice:
\[
s_{xx} = 0 \implies s_x = C_1 \implies s = C_1 x + C_2
\]
Applying the boundary conditions, we get
\[
s_x(0) = C_1 = 1 \implies C_1 = 1
\]
\[
s_x(1) = C_1 = 1 \implies C_1 = 1
\]
Thus, the steady-state solution is
\[
s(x) = x + C_2
\]
\textbf{Transient Solution}
Since the solution to the entire problem is $u(x,t) = s(x) + v(x,t)$, we can substitute this into the original PDE to find the transient part of the solution:
\[
u(x,t) = s(x) + v(x,t)
\]
We need to find $v(x,t)$ so:
\[
v(x,t) = u(x,t) - s(x) = u(x,t) - (x + C_2)
\]
This makes the new PDE for $v$:
\[
\begin{cases}
v_t = v_{xx}, & 0 < x < 1,\ t > 0, \\
v_x(0,t) = v_x(1,t) = 0, & t > 0, \\
v(x,0) = u(x,0) - s(x) = \cos^2(\pi x) - C_2, & 0 \leq x \leq 1.
\end{cases}
\]
This is a homogeneous heat equation with homogeneous Neumann boundary conditions. We can solve this using separation of variables or by inspection. Let $v(x,t) = X(x)T(t)$, then we have
\[
\frac{T'}{T} = \frac{X''}{X} = -\lambda^2
\]
\[
\implies T(t) = e^{-\lambda^2 t}
\]
\[
X'' + \lambda^2 X = 0
\implies X(x) = A \cos(\lambda x) + B \sin(\lambda x)
\]
Checking the boundary conditions, we get
\[
X'(0) = -A \lambda \sin(0) + B \lambda \cos(0) = B \lambda = 0 \implies B = 0
\]
\[
X'(1) = -A \lambda \sin(\lambda) = 0 \implies \lambda = n\pi, \ n = 0,1,2,\dots
\]
Thus, the transient solution is
\[
v(x,t) = \frac{A_0}{2} + \sum_{n=1}^{\infty} A_n \cos(n\pi x) e^{-(n\pi)^2 t}
\]
To find the coefficients $A_n$, we use the initial condition for $v$:
\[
v(x,0) = \cos^2(\pi x) - C_2 = \frac{A_0}{2} + \sum_{n=1}^{\infty} A_n \cos(n\pi x)
\]
\[
\implies \frac{1}{2} + \frac{1}{2} \cos(2\pi x) - C_2 = \frac{A_0}{2} + \sum_{n=1}^{\infty} A_n \cos(n\pi x)
\]
Let $C_2 = \frac{1}{2}$, then $A_0 = 0$ and $A_2 = \frac{1}{2}$. Thus, the transient solution is
\[
v(x,t) = \frac{1}{2} \cos(2\pi x) e^{-4\pi^2 t}
\]
\textbf{Final Solution}
Thus, the final solution to the problem is
\[
\boxed{
u(x,t) = s(x) + v(x,t) = x + \frac{1}{2} + \frac{1}{2} \cos(2\pi x) e^{-4\pi^2 t}
}
\]

\prob{3}
Solve the heat-flow problem
\[
\begin{cases}
u_t = u_{xx} - 0.1\,u, & 0 < x < \pi,\ t > 0, \\
u(0,t) = u_x(\pi,t) = 0, & t > 0, \\
u(x,0) = \sin(1.5\,x), & 0 \leq x \leq \pi.
\end{cases}
\]

\sol
This is a homogeneous heat equation with a reaction term. We can solve this using separation of variables. Let $u(x,t) = X(x)T(t)$, then we have
\[
X(x) T'(t) = X''(x) T(t) - 0.1 X(x) T(t)
\implies X(x) T'(t) + 0.1 X(x) T(t) = X''(x) T(t)
\]
\[
\implies \frac{T'}{T} + 0.1 = \frac{X''}{X} = -\lambda^2
\]
\[
\implies T(t) = e^{-(\lambda^2 + 0.1) t}
\]
\[
X'' + \lambda^2 X = 0
\implies X(x) = A \cos(\lambda x) + B \sin(\lambda x)
\]
Applying the boundary conditions, we get
\[
X(0) = A = 0 \implies A = 0
\]
\[
X'(\pi) = B \lambda \cos(\lambda \pi) = 0 \implies \lambda = \frac{n}{2}, \ n = 1,3,5,\dots
\]
Thus, the solution is
\[
u(x,t) = \sum_{n=1}^\infty B_n \sin\left(\frac{n}{2} x\right) e^{-\left(\left(\frac{n}{2}\right)^2 + 0.1\right) t}
\]
To find the coefficients $B_n$, we use the initial condition:
\[
u(x,0) = \sin(1.5 x) = \sum_{n=1}^\infty B_n \sin\left(\frac{n}{2} x\right)
\]
The only term in the series that will be nonzero is when $n=3$ because $\sin(1.5 x) = \sin\left(\frac{3}{2} x\right)$. Thus, we have
\[
B_3 = 1
\]
and all other $B_n$ are zero. Thus, the solution is
\[
\boxed{
u(x,t) = \sin\left(\frac{3}{2} x\right) e^{-\left(\left(\frac{3}{2}\right)^2 + 0.1\right) t}
}
\]

\prob{4}
Solve the diffusion-convection problem
\[
\begin{cases}
u_t = u_{xx} - 2u_x, & 0 < x < \pi,\ t > 0, \\
u(0,t) = u(\pi,t) = 0, & t > 0, \\
u(x,0) = e^x \sin(x), & 0 \leq x \leq \pi,
\end{cases}
\]
by first transforming it into an easier problem.

\sol
This is a diffusion-convection problem. There are two ways to solve this problem: use the same trick in the previous problem to transform it into a homogeneous heat equation, or use straight separation of variables. \\
\textbf{Method 1: Separation of Variables}\\
I like separation of variables better since I am stronger at ODEs. Let $u(x,t) = X(x)T(t)$, then we have
\[
X(x) T'(t) = X''(x) T(t) - 2 X'(x) T(t) = \tup{X''(x) -2X'(x)}T(t)
\]
\[
\implies
\frac{T'}{T} = \frac{X'' - 2X'}{X} = -\lambda^2
\]
\[
\implies T(t) = e^{-\lambda^2 t}
\]
\[
X'' - 2X' + \lambda^2 X = 0
\]
This is a second-order linear ODE with constant coefficients. The characteristic equation is
\[
r^2 - 2r + \lambda^2 = 0
\implies r = \frac{2 \pm \sqrt{4-4\lambda^2}}{2} = 1 \pm \sqrt{1-\lambda^2}
\]
Thus, the solution to the ODE is
\[
X(x) = e^x \left(A \cos\left(\sqrt{\lambda^2 - 1} x\right) + B \sin\left(\sqrt{\lambda^2 - 1} x\right)\right)
\]
Applying the boundary conditions, we get
\[
X(0) = e^0 (A \cos(0) + B \sin(0)) = A = 0 \implies A = 0
\]
\[
X(\pi) = e^\pi B \sin\left(\sqrt{\lambda^2 - 1} \pi\right) = 0 
\]
\[
\implies \sin \tup{\sqrt{\lambda^2-1}\pi} = 0 \implies \sqrt{\lambda^2 - 1} = n, \ n = 1,2,3,\dots
\]
\[
\implies \lambda^2 = n^2 + 1, \ n = 1,2,3,\dots
\]
Thus, the solution is
\[
u(x,t) = \sum_{n=1}^\infty B_n e^x \sin(nx) e^{-(n^2 + 1) t}
\]
To find the coefficients $B_n$, we use the initial condition:
\[
u(x,0) = e^x \sin(x) = \sum_{n=1}^\infty B_n e^x \sin(nx)
\]
The only term in the series that will be nonzero is when $n=1$ because $e^x \sin(x) = e^x \sin(1 \cdot x)$. Thus, we have
\[
B_1 = 1
\]
and all other $B_n$ are zero. Thus, the solution is
\[
\boxed{
u(x,t) = e^x \sin(x) e^{-2t}
}
\]
\textbf{Method Two: Converting it into an easier problem}\\
Let $v(x,t) = e^{-x} u(x,t)$, then we have
\[
\begin{cases}
v_t = e^{-x} u_t = e^{-x} (u_{xx} - 2u_x) = e^{-x} u_{xx} - 2 e^{-x} u_x, & 0 < x < \pi,\ t > 0, \\
v(0,t) = e^0 u(0,t) = 0,\quad v(\pi,t) = e^{-\pi} u(\pi,t) = 0, & t > 0, \\
v(x,0) = e^{-x} u(x,0) = e^{-x} e^x \sin(x) = \sin(x), & 0 \leq x \leq \pi,
\end{cases}
\]
which is a homogeneous heat equation with homogeneous Dirichlet boundary conditions. We can solve this using separation of variables or by inspection. Let $v(x,t) = X(x)T(t)$, then we have
\[
\frac{T'}{T} = \frac{X''}{X} = -\lambda^2
\]
\[
\implies T(t) = e^{-\lambda^2 t}
\]
\[
X'' + \lambda^2 X = 0
\implies X(x) = A \cos(\lambda x) + B \sin(\lambda x)
\]
Applying the boundary conditions, we get
\[
X(0) = A = 0 \implies A = 0
\]
\[
X(\pi) = B \sin(\lambda \pi) = 0 \implies \lambda = n, \ n = 1,2,3,\dots
\]
Thus, the solution is
\[
v(x,t) = \sum_{n=1}^\infty B_n \sin(nx) e^{-n^2 t}
\]
To find the coefficients $B_n$, we use the initial condition:
\[
v(x,0) = \sin(x) = \sum_{n=1}^\infty B_n \sin(nx)
\]
The only term in the series that will be nonzero is when $n=1$ because $\sin(x) = \sin(1 \cdot x)$. Thus, we have
\[
B_1 = 1
\]
and all other $B_n$ are zero. Thus, the solution is
\[
v(x,t) = \sin(x) e^{-t}
\]
Finally, we can convert back to $u$:
\[
\boxed{
u(x,t) = e^x v(x,t) = e^x \sin(x) e^{-t} = e^x \sin(x) e^{-2t}
}
\]


\prob{5}
Solve the initial-boundary value problem with a source term
\[
\begin{cases}
u_t = u_{xx} + \sin(3\pi x), & 0 < x < 1,\ t > 0, \\
u(0,t) = u(1,t) = 0, & t > 0, \\
u(x,0) = \sin(\pi x), & 0 \leq x \leq 1.
\end{cases}
\]
Identify the steady-state and transient parts of the solution.

\sol
This problem is nonhomogeneous, so the solution will contain a steady-state part and a transient part. The solution looks like the following:
\[
u(x,t) = s(x) + v(x,t),
\]
where $s(x)$ is the steady-state solution and $v(x,t)$ is the transient part of the solution. The steady-state solution is the solution to the following boundary value problem because the transient part will vanish as $t \to \infty$:
\[
\begin{cases}
s_{xx} + \sin(3\pi x) = 0, & 0 < x < 1, \\
s(0) = 0,\quad s(1) = 0.
\end{cases}
\]
\textbf{Steady-State Solution}
To solve this, we integrate the equation twice:
\[
s_{xx} + \sin(3\pi x) = 0 \implies s_{xx} = -\sin(3\pi x)
\]
\[
s_x = \int -\sin(3\pi x) dx = \frac{1}{3\pi} \cos(3\pi x) + C_1
\]
\[
s = \int \left(\frac{1}{3\pi} \cos(3\pi x) + C_1\right) dx = \frac{1}{9\pi^2} \sin(3\pi x) + C_1 x + C_2
\]
Applying the boundary conditions, we get
\[
s(0) = C_2 = 0 \implies C_2 = 0
\]
\[
s(1) = \frac{1}{9\pi^2} \sin(3\pi) + C_1 = 0 \implies C_1 = 0
\]
Thus, the steady-state solution is
\[
s(x) = \frac{1}{9\pi^2} \sin(3\pi x)
\]
\textbf{Transient Solution}
Since the solution to the entire problem is $u(x,t) = s(x) + v(x,t)$, we can substitute this into the original PDE to find the transient part of the solution:
\[
u(x,t) = s(x) + v(x,t)
\]
We need to find $v(x,t)$ so:
\[
v(x,t) = u(x,t) - s(x) = u(x,t) - \frac{1}{9\pi^2} \sin(3\pi x)
\]
This makes the new PDE for $v$:
\[
\begin{cases}
v_t = v_{xx}, & 0 < x < 1,\ t > 0, \\
v(0,t) = 0,\quad v(1,t) = 0, & t > 0, \\
v(x,0) = u(x,0) - s(x) = \sin(\pi x) - \frac{1}{9\pi^2} \sin(3\pi x), & 0 \leq x \leq 1.
\end{cases}
\]
This is a homogeneous heat equation with homogeneous Dirichlet boundary conditions. We can solve this using separation of variables or by inspection. Let $v(x,t) = X(x)T(t)$, then we have
\[
\frac{T'}{T} = \frac{X''}{X} = -\lambda^2
\]
\[
\implies T(t) = e^{-\lambda^2 t}
\]
\[
X'' + \lambda^2 X = 0
\implies X(x) = A \cos(\lambda x) + B \sin(\lambda x)
\]
Applying the boundary conditions, we get
\[
X(0) = A = 0 \implies A = 0
\]
\[
X(1) = B \sin(\lambda) = 0 \implies \lambda = n\pi, \ n = 1,2,3,\dots
\]
Thus, the transient solution is
\[
v(x,t) = \sum_{n=1}^{\infty} B_n \sin(n\pi x) e^{-(n\pi)^2 t}
\]
To find the coefficients $B_n$, we use the initial condition for $v$:
\[
v(x,0) = \sin(\pi x) - \frac{1}{9\pi^2} \sin(3\pi x) = \sum_{n=1}^{\infty} B_n \sin(n\pi x)
\]
The only terms in the series that will be nonzero are when $n=1$ and $n=3$ because $\sin(\pi x) = \sin(1 \cdot \pi x)$ and $\sin(3\pi x) = \sin(3 \cdot \pi x)$. Thus, we have
\[
B_1 = 1, \quad B_3 = -\frac{1}{9\pi^2}
\]
and all other $B_n$ are zero. Thus, the transient solution is
\[
v(x,t) = \sin(\pi x) e^{-\pi^2 t} - \frac{1}{9\pi^2} \sin(3\pi x) e^{-9\pi^2 t}
\]
\textbf{Final Solution}
Thus, the final solution to the problem is
\[
\boxed{
u(x,t) = s(x) + v(x,t) = \frac{1}{9\pi^2} \sin(3\pi x) + \sin(\pi x) e^{-\pi^2 t} - \frac{1}{9\pi^2} \sin(3\pi x) e^{-9\pi^2 t}
}
\]
\end{document}